\subsection{EXTENSION OF RATIONAL FUNCTIONS}

If the results of no. 8 are applied to the vector space \(\mathbf{V} = \mathbf{K}_s\) of dimension 1, it is seen that there exists a canonical injection \(\phi\) of \(\mathbf{K}_s\) into the projective line \(\mathbf{P}_1(\mathbf{K}) = \mathbf{P}(\mathbf{K}_s \times \mathbf{K}_s)\) ; for all \(\xi \in \mathbf{K}\) , \(\phi(\xi)\) is the point with homogeneous coordinates \((1, \xi)\) with respect to the canonical basis ( \(\S 1\) , no. 11) of \(\mathbf{K}_s \times \mathbf{K}_s\) . The complement of \(\phi(\mathbf{K})\) in \(\mathbf{P}_1(\mathbf{K})\) consists of the single point with homogeneous coordinates \((0, 1)\) with respect to the above basis; it is called the ``point at infinity''. \(\mathbf{P}_1(\mathbf{K})\) is also called the projective field associated with \(\mathbf{K}\) and denoted by \(\tilde{\mathbf{K}}\) , the point at infinity in \(\tilde{\mathbf{K}}\) being denoted by \(\infty\) .

Consider in particular the case where K is a commutative field and let \(f \in \mathbf{K}(\mathbf{X})\) be a rational function in one indeterminate over \(\mathbf{K}\) (IV, §4); if \(f \neq 0\) , there is a unique expression \(f = \alpha p/q\) , where \(\alpha \in \mathbf{K}^*\) and \(p\) and \(q\) are two relatively prime monic polynomials (VII, §1); let \(m\) and \(n\) be their respective degrees and let \(r = \sup(m, n)\) . We write

\[
p _ {1} (\mathrm{T}, \mathrm{X}) = \mathrm{T} ^ {r} p (\mathrm{X} / \mathrm{T}), \quad q _ {1} (\mathrm{T}, \mathrm{X}) = \mathrm{T} ^ {r} q (\mathrm{X} / \mathrm{T});
\]

\(p_{1}\) and \(q_{1}\) are two homogeneous polynomials of degree r over K such that \(p(\mathbf{X}) = p_{1}(1, \mathbf{X})\) , \(q(\mathbf{X}) = q_{1}(1, \mathbf{X})\) . Hence, for every element \(\xi \in K\) which is not a zero of \(q(\mathbf{X})\) , \(f(\xi) = \alpha p(\xi)/q(\xi)\) is defined and we may write

\[
f (\xi) = \alpha p _ {1} (1, \xi) / q _ {1} (1, \xi) = \alpha p _ {1} (\lambda , \lambda \xi) / q _ {1} (\lambda , \lambda \xi)
\]

for all \(\lambda \neq 0\) in K. Consider then the mapping

\[
(\eta , \xi) \mapsto (q _ {1} (\eta , \xi), p _ {1} (\eta , \xi))
\]

of \(K^{2}\) into itself; it is compatible with the equivalence relation \(\Delta(\mathbf{K}^{2})\) and therefore defines, when passing to the quotients, a mapping \(\tilde{f}\) of \(\tilde{K}\) into itself which coincides with \(\xi\mapsto f(\xi)\) at the points where this rational function is defined; it is said, by an abuse of language, that \(\tilde{f}\) is the canonical extension of f to \(\tilde{K}\) .

For example, if \(f = 1 / \mathbf{X}\) , then \(\tilde{f}(0) = \infty\) and \(\tilde{f}(\infty) = 0\) ; if

\[
f = (a \mathbf {X} + b) / (c \mathbf {X} + d)
\]

with \(ad - bc \neq 0\) , then \(\tilde{f}(-d/c) = \infty\) , \(\tilde{f}(\infty) = a/c\) if \(c \neq 0\) , \(\tilde{f}(\infty) = \infty\) if \(c = 0\) . If \(f = a_0\mathbf{X}^n + \cdots + a_n\) is a polynomial of degree \(n > 0\) , then \(\tilde{f}(\infty) = \infty\) .

